\section{Evaluation}
\label{sec:evaluation}

All reported $p$-values are exact, two-sided Mann-Whitney U tests
(\texttt{scipy.stats.mannwhitneyu}, \texttt{method="exact"}), computed
directly rather than via the closed-form shortcut a fixed-budget design can
tempt one into approximating by hand; at $n{=}5$ per arm with complete rank
separation, $C(10,5) = 252$ and exactly two of those partitions produce
complete separation in either direction, so $p = 2/252 \approx 0.00794$ is
the exact value in every complete-separation case reported below, not a
coincidence between unrelated effects.

\subsection{RQ1: does the refinement mechanism transfer?}
\label{sec:evaluation-rq1}
Table~\ref{tab:rq1} reports acceptance rate -- the fraction of generated
documents that are syntactically valid JSON by the same criterion the
original grammar uses, any value type at the top level -- for five seeded
baseline runs and five seeded refined runs. The baseline is exactly 50.0\%
on every run by construction (its input stream deterministically alternates
a grammar-valid generator with a deliberately corrupted variant of it).
Every refined run exceeds every baseline run: complete rank separation,
exact $p \approx 0.00794$. One refined run's lower rate (68.6\%) traces to a
seed on which only one of five iterations produced usable data, the rest
consumed by ordinary LLM code-generation failures (e.g.\ one proposal built
an unbounded-depth strategy that hit Python's recursion limit) -- logged
and handled exactly as the original pipeline's error-recovery design
specifies, not a defect in this port.

\begin{table}[t]
\centering
\footnotesize
\caption{RQ1: acceptance rate, $n{=}5$ seeded runs per arm.}
\label{tab:rq1}
\begin{tabularx}{\columnwidth}{@{}lX@{}}
\toprule
Metric & Value \\
\midrule
Baseline, per-run (\%) & 50.0, 50.0, 50.0, 50.0, 50.0 \\
Refined, per-run (\%) & 68.6, 97.3, 98.0, 98.1, 98.2 \\
Baseline mean (stdev) & 50.0 (0.0) \\
Refined mean (stdev) & 92.0 (13.1) \\
\bottomrule
\end{tabularx}
\end{table}

This reproduces the original study's cJSON result in shape -- 58.2\% to
97.1\% there, 50.0\% to a 92.0\% mean here -- on a completely different
runtime and language, with the identical refinement-loop mechanism and no
target-specific changes to the loop itself. We read this as confirmation
that the mechanism, not merely the original numeric result, generalizes.

\subsection{RQ2: does refinement find more divergence than a static generator?}
\label{sec:evaluation-rq2}
The first refinement prompt (Section~\ref{sec:rq2-prompt}) was run across
three full seeded runs (seeds 0-2) before any revision: zero cross-path
divergence in every one of the resulting iteration-campaigns, despite
per-path rejection counts climbing sharply by the final iteration of each
run (up to 208 of 264 schema-evaluated documents on one run). This ruled
out single-run variance as an explanation before the prompt was changed.
The revised prompt (Section~\ref{sec:rq2-prompt}) was validated on one
held-out seed (11.8\% divergence in its best iteration, against a confirmed
0\% baseline under the old prompt) before being adopted for the five-seed
comparison reported in Table~\ref{tab:rq2}.

\begin{table}[t]
\centering
\footnotesize
\caption{RQ2: cross-path divergence rate, $n{=}5$ seeded runs per arm. Baseline uses the static schema-aware generator; refined uses the corrected prompt.}
\label{tab:rq2}
\begin{tabularx}{\columnwidth}{@{}lX@{}}
\toprule
Metric & Value \\
\midrule
Baseline, per-run (\%) & 19.8, 21.8, 22.2, 23.4, 23.8 \\
Refined, per-run (\%) & 1.3, 1.45, 2.75, 3.6, 5.53 \\
Baseline mean (stdev) & 22.2 (1.57) \\
Refined mean (stdev) & 2.93 (1.74) \\
\bottomrule
\end{tabularx}
\end{table}

Every baseline run exceeds every refined run: complete rank separation,
exact $p \approx 0.00794$, in the opposite direction from RQ1. Restricting
the comparison to only the subset of each run's generated documents that
reached the schema stage at all (i.e.\ excluding the refined arm's
additional overhead from syntactically invalid JSON, since the sandbox
forbids importing a JSON library and the LLM must hand-build JSON text)
still leaves a wide gap: a 4.10\% mean divergence rate among
schema-evaluated documents for the refined arm, against the baseline's
22.2\% (identical to its overall rate, since the static generator always
emits valid JSON). Generation-quality overhead is real -- the refined
arm's schema-evaluated fraction of all generated documents ranged
41.9\%-83.7\% across runs, versus the baseline's fixed 100\% -- but this
comparison alone cannot separate how much of the gap it explains from a
targeting shortfall; Section~\ref{sec:evaluation-ablation} isolates the two
directly. Pooling every schema-evaluated record across all
five refined seeds, the loop found exactly the same four divergence
patterns the static generator found (including both confirmed defects
below), never a qualitatively new one, at lower absolute rate throughout.

\subsection{Ablation: generation overhead, or targeting?}
\label{sec:evaluation-ablation}
Table~\ref{tab:ablation} reports the ablation from Section~\ref{sec:ablation-json}:
an otherwise-identical refined arm, five fresh seeds (100-104, disjoint
from every other reported run), with \texttt{import json} additionally
permitted. Every one of the 24 iterations that produced data reached
100\% (or, in one case affected by an unrelated encoding issue, 82.5\%)
\texttt{schema\_evaluated} -- the generation-overhead problem is
essentially eliminated, exactly as expected once the LLM no longer needs
to hand-roll JSON text.

\begin{table}[t]
\centering
\footnotesize
\caption{Ablation: divergence rate with \texttt{import json} permitted, $n{=}5$ fresh seeded runs, vs.\ the main text's constrained-sandbox refined arm and static baseline.}
\label{tab:ablation}
\begin{tabularx}{\columnwidth}{@{}lX@{}}
\toprule
Metric & Value \\
\midrule
\texttt{json}-allowed, per-run (\%) & 0.05, 2.2, 2.36, 2.36, 4.05 \\
\texttt{json}-allowed mean (stdev) & 2.20 (1.42) \\
Constrained-sandbox mean (stdev) & 2.93 (1.74) \\
Static baseline mean (stdev) & 22.2 (1.57) \\
\bottomrule
\end{tabularx}
\end{table}

The divergence rate did not improve: 2.20\% mean, statistically
indistinguishable from the constrained sandbox's 2.93\% (exact two-sided
Mann-Whitney U, $p \approx 0.69$), and both remain in complete separation
below the static baseline's 22.2\% ($p \approx 0.00794$ against the
\texttt{json}-allowed arm too). This is a clean, decisive answer to the
question Section~\ref{sec:evaluation-rq2} could only partially resolve by
conditioning on schema-evaluated documents: removing the sandbox's
generation-overhead entirely, verified directly rather than inferred, does
not close any measurable part of the RQ2 gap. The shortfall is a targeting
problem, not a JSON-syntax-generation problem.

\subsection{Follow-up: does richer feedback close the targeting gap?}
\label{sec:evaluation-category-feedback}
Section~\ref{sec:evaluation-ablation} established the gap is a targeting
problem; the direct follow-up question is whether telling the proposer
\emph{which categories} of perturbation its divergence-causing documents
last exhibited -- missing field, wrong type, null override, unrecognized
enum, boundary-magnitude \texttt{id}, extra key, deep nesting, computed
automatically from already-collected data and fed back without naming
specific schema fields -- helps it target more effectively than the single
numeric score alone. We ran a third $n{=}5$ arm (fresh seeds 200-204) under
this richer prompt, otherwise identical to the main text's constrained
sandbox.

\begin{table}[t]
\centering
\footnotesize
\caption{Follow-up: divergence rate with perturbation-category feedback, $n{=}5$ fresh seeded runs, vs.\ the constrained-sandbox refined arm.}
\label{tab:category-feedback}
\begin{tabularx}{\columnwidth}{@{}lX@{}}
\toprule
Metric & Value \\
\midrule
Category-feedback, per-run (\%) & 0.0, 0.0, 1.0, 1.2, 2.55 \\
Category-feedback mean (stdev) & 0.95 (1.05) \\
Constrained-sandbox (score-only) mean (stdev) & 2.93 (1.74) \\
\bottomrule
\end{tabularx}
\end{table}

Richer feedback did not help -- it did measurably worse: a 0.95\% mean
divergence rate against the score-only prompt's 2.93\%, a difference that
holds under an exact two-sided Mann-Whitney U both over all generated
documents ($p \approx 0.032$) and restricted to schema-evaluated documents
only ($p \approx 0.032$, ruling out proposal-failure noise as the
explanation). Two of the five category-feedback seeds found literally zero
divergence in their only successfully-validated iteration, not merely zero
because every iteration failed to produce data. We do not have a confirmed
account of why more detailed feedback underperforms simpler feedback here;
a plausible but untested reading is that a longer prompt enumerating seven
named categories and their definitions competes for the proposer's
attention with the actual code-generation task, or invites the model to
write more elaborate, more bug-prone generator code attempting to
explicitly target each named category, rather than the smaller, more
robust generators the shorter prompt tends to produce. This result is
small-sample ($n{=}5$) and reported descriptively for that reason, but it
is a genuine data point against the intuitive assumption that more
informative feedback is strictly better for this task.

\subsection{RQ3: does numeric decoding preserve exact values?}
\label{sec:evaluation-rq3}
Table~\ref{tab:rq3} reports the ground-truth \texttt{id} comparison
(Section~\ref{sec:oracles}) computed as an analysis pass over
schema-evaluated records already collected across every RQ2 campaign in
this paper (static baseline, constrained-sandbox refined arm, and the
Section~\ref{sec:evaluation-ablation} ablation, pooled) -- no separate
generator, harness, or campaign was needed. The schema's recursive
\texttt{child} field means an accepting path's canonical value mirrors the
full decoded tree, not just its top-level fields, so every nesting depth
already present in collected data is another independent check of the same
kind; walking all of them (rather than only the top level, as an earlier
pass over a smaller pool of campaigns did) raises $N$ to 43{,}336 checks.

\begin{table}[t]
\centering
\footnotesize
\caption{RQ3: exact-integer round-trip on the \texttt{id} field, $N{=}43{,}336$ checks (all nesting depths, pooled across every RQ2 campaign).}
\label{tab:rq3}
\begin{tabular}{@{}lrrr@{}}
\toprule
Path & Accepted & Mismatched & Rate \\
\midrule
Manual & 10{,}070 & 0 & 0.0\% \\
json\_serializable & 11{,}449 & 649 & 5.7\% \\
freezed & 11{,}449 & 649 & 5.7\% \\
built\_value & 10{,}368 & 0 & 0.0\% \\
\bottomrule
\end{tabular}
\end{table}

Manual parsing and \texttt{built\_value} never silently mismatch: across
ten thousand accepted decodes each, they either decode \texttt{id} exactly
or reject the input outright. \texttt{json\_serializable} and
\texttt{freezed} are identical to each other (freezed delegates this
field's decoding to json\_serializable's generated code) and both show a
5.7\% silent mismatch rate overall. Every observed mismatch saturates to
exactly \texttt{int64::MAX} or \texttt{int64::MIN} -- e.g.\ an input
\texttt{id} of $-9{,}223{,}372{,}036{,}854{,}775{,}810$ decodes to
$-9{,}223{,}372{,}036{,}854{,}775{,}808$ -- consistent with the generated
code calling \texttt{(json['id'] as num).toInt()} rather than
\texttt{json['id'] as int}: an out-of-range JSON integer literal is first
promoted to a \texttt{double} by \texttt{dart:convert}, and \texttt{toInt()}
on an out-of-range \texttt{double} saturates silently rather than throwing.

Breaking \texttt{json\_serializable}'s rate out by nesting depth shows it
is not uniform: 3.4\% at depth 0 ($n{=}8{,}198$), rising to 7.7\% at depth 1
($n{=}2{,}377$), 20.7\% at depth 2 ($n{=}590$), and 21.2\% at depth 3
($n{=}240$), before the sample thins out past depth 3 ($n \le 23$ per level,
too small to read as more than suggestive). We do not have a confirmed
mechanistic account of why deeper nesting correlates with a higher
mismatch rate -- a plausible but untested guess is that the generator's
perturbation and recursion machinery are not independent, so documents that
happen to nest deeply also happen to disproportionately carry
boundary-magnitude \texttt{id} values -- and report the pattern
descriptively rather than as a tested effect of depth itself.

\subsection{A targeted side study: does \texttt{freezed} ever diverge from \texttt{json\_serializable}?}
\label{sec:evaluation-union}
In every result reported above, \texttt{freezed} and \texttt{json\_serializable}
are identical, because \texttt{freezed} delegates every Record field's
leaf-decoding to json\_serializable's generated code
(Section~\ref{sec:target-set}) -- the schema never exercises the one
feature (native discriminated-union JSON dispatch) that is actually
\texttt{freezed}'s own. We ran a small, separate, deliberately qualitative
side study -- not folded into the RQ1-RQ3 statistical results, and not
claiming their level of rigor -- to check whether this identity holds even
where it should not have to: a two-variant discriminated union,
\texttt{Payload.text(String)} / \texttt{Payload.number(num)}, discriminated
by a \texttt{"type"} key, implemented once per path following each
framework's own idiomatic pattern (manual: an if/else dispatch;
json\_serializable: a hand-written dispatch function plus
\texttt{@JsonSerializable()} leaf classes, since json\_serializable has no
native polymorphism; freezed: \texttt{@Freezed(unionKey: 'type')}, its own
built-in mechanism; built\_value: a hand-written dispatch function plus a
custom \texttt{Serializer} per leaf, since built\_value has no native
polymorphism either). Thirteen hand-picked documents (both valid variants,
an unknown/missing/null/wrong-type/differently-cased discriminant, and
wrong-type or missing/null leaf values) were run through all four paths.

The result: \texttt{freezed} finally diverges from \texttt{json\_serializable}
-- but at the exception-\emph{type} level, not accept/reject or value.
Every invalid-discriminant case (unknown, missing, \texttt{null},
wrong-type, or differently-cased) is rejected by all four paths, but
\texttt{freezed} raises \texttt{CheckedFromJsonException} (a public,
structured, field-attributed error type from \texttt{json\_annotation}'s
own checked-decoding support) while manual, json\_serializable, and
built\_value all raise a generic \texttt{ArgumentError} -- the first
behavioral difference of any kind observed between B and C anywhere in
this paper. For wrong-type or missing leaf-\emph{value} cases (the
discriminant itself is valid), all four paths converge back to the same
pattern as the main schema: A/B/C raise the private \texttt{\_TypeError},
D raises its own \texttt{DeserializationError}, and no B-vs-C difference
appears. No accept/reject divergence and no value divergence (all four
accepting but disagreeing on the decoded value) were observed in this
small sample; the discriminant-based dispatch in all four paths is simple
and deterministic enough that we would not expect the latter without a
schema feature (e.g.\ ambiguous or overlapping variant matching) this
minimal union does not have. We report this as a real, novel,
practically-relevant qualitative finding -- json\_serializable's manual
dispatch pattern for polymorphic JSON gives callers a less specific
exception to catch than freezed's native union support does, for the exact
failure mode (a malformed discriminant) polymorphic JSON is most likely to
hit in practice -- rather than as a quantitative result on the level of
Sections~\ref{sec:evaluation-rq1}--\ref{sec:evaluation-category-feedback}:
formalizing it into an automated oracle, generator, and seeded campaign at
scale is future work (Section~\ref{sec:future-work}), not attempted here.

\subsection{Two concrete, practitioner-actionable defects}
\label{sec:evaluation-defects}
Both of the following were found by the static generator alone, with no
LLM involved, and are independent of the RQ1/RQ2 statistical comparisons
above. \textbf{(1) Silent integer saturation}, RQ3 above: any Flutter
application using \texttt{json\_serializable} or \texttt{freezed} to
decode an integer identifier (a snowflake ID, an external database key)
that can exceed 64-bit range will silently receive
\texttt{9223372036854775807} rather than an exception, for any such
identifier, with no warning. \textbf{(2) Silent default on a missing
collection field}: \texttt{built\_value} decodes a document with the
\texttt{tags} field entirely absent by silently defaulting it to an empty
list, while manual parsing, \texttt{json\_serializable}, and
\texttt{freezed} all correctly reject the same input as malformed. This
does not generalize to missing scalar fields -- \texttt{id},
\texttt{amount}, and \texttt{status} are all correctly rejected as missing
by every path including \texttt{built\_value} -- so the finding is
precisely scoped to \texttt{built\_value}'s list-typed field construction,
not general laxness. Both defects are reproducible with the released
harness and require no fuzzing infrastructure to confirm by hand.
